\documentclass[a4paper,10pt]{ltjsarticle}
\usepackage[margin=25mm]{geometry}
\usepackage{amsmath,amssymb,amsthm,mathtools,bm}
\usepackage{enumitem,booktabs,array,longtable}
\usepackage[unicode,hidelinks]{hyperref}
\usepackage{bookmark}
\setlength{\emergencystretch}{2em}
\newtheorem{theorem}{定理}[section]
\newtheorem{proposition}[theorem]{命題}
\newtheorem{lemma}[theorem]{補題}
\newtheorem{corollary}[theorem]{系}
\theoremstyle{definition}
\newtheorem{definition}[theorem]{定義}
\newtheorem{hypothesis}[theorem]{仮定}
\theoremstyle{remark}
\newtheorem{remark}[theorem]{注意}
\newcommand{\Perf}{\operatorname{Perf}}
\newcommand{\RHom}{\operatorname{RHom}}
\newcommand{\REnd}{\operatorname{REnd}}
\newcommand{\thick}{\operatorname{thick}}
\newcommand{\Cone}{\operatorname{Cone}}
\newcommand{\Ku}{\operatorname{Ku}}
\newcommand{\qgr}{\operatorname{qgr}}
\newcommand{\QGr}{\operatorname{QGr}}
\newcommand{\gr}{\operatorname{gr}}
\newcommand{\Dsg}{D_{\mathrm{sg}}^b}
\newcommand{\CC}{\mathbb C}
\newcommand{\OO}{\mathcal O}
\newcommand{\cC}{\mathcal C}
\newcommand{\cK}{\mathcal K}
\newcommand{\cD}{\mathcal D}
\newcommand{\otL}{\otimes^{\mathbf L}}
\newcommand{\dd}{^{\vee}}
\newcommand{\LL}{\mathbf L}
\newcommand{\RR}{\mathbf R}
\title{非可換三次超曲面の局所rotationと\protect\\可換復元のsaturation還元}
\author{nc-cubic 研究原稿C・独立読解用の現行節}
\date{2026年11月10日・v21}
\begin{document}
\maketitle
\begin{abstract}
滑らかな三次超曲面に特殊化するproperな局所dg模型から出発し、標準rotation
$\rho=\LL_{\OO}\circ(1)$の逆関手を$(\!-1)\circ\RR_{\OO}$として構成する。
証明は一つのSerre traceを持ち上げるものであり、残余dg圏のsmoothnessを必要としない。
さらに残余の対角がperfectなら、線束生成部分と全coherent導来圏の同定、
およびrotation付きpairの有限型開集合への降下が従うことを示す。
さらに同値双加群の持上げ障害をrotation不変部分で消去し、
実在するsmooth properなdg族に対するétale局所実現原理を証明する。
BFKの固定sector公式から全次元のrotation不変な三次Hochschild障害の消滅を示し、
特に非楕円cubic fivefoldで\(\mathrm{HH}^3=0\)を得る。
従って、所定の局所量子化が得られたPoisson成分のgenericな可換復元は、
残余対角のperfect性へ還元される。一般の残余smoothnessは未証明である。
\end{abstract}
\tableofcontents

\section{問題、既知の入力、主結果}\label{sec:intro}
本研究の目標は、任意の$n\ge3$と滑らかな複素三次超曲面$X^n$について、
そのPoissonスキームの各被約既約成分の幾何学的一般点で、構成した量子化の残余dg圏を
可換な滑らかな三次超曲面のrotation付きKuznetsov成分として実現することである。
全てのAS-regular中心三次商の全点に同じ存在を要求する問題ではない。

以下の全ての圏、関手、自然変換にはdg enhancementを保ち、底変更はMorita底変更を意味する。
圏の同値を単なる三角圏の同値へ弱めない。
$R$上properとは対象間の全導来Homが$R$-perfectであること、smoothとは対角が
包絡dg圏上perfectであることをいう。
$V^\vee=\RHom_R(V,R)$とし、直交条件は全cohomological shiftで課す。

\begin{hypothesis}[局所圏論的入力]\label{hyp:categorical}
$(R,\mathfrak m,k)$はnoetherian局所環である。
$\cC$は$R$-properな冪等完備pretriangulated dg圏で、dg自己同値$T$と対象$\OO$を持つ。
$\OO_i=T^i\OO$と書き、有限個の$\OO_i$の直和$E$が$\cC$を古典的に生成する。
ある滑らかな三次超曲面$X\subset\mathbb P^{n+1}_k$、$n\ge3$に対し、指定された同値
\[
 (\cC_k,T_k,\OO_k)\simeq(\Perf(X),-\otimes\OO_X(1),\OO_X)
 \tag{L}\label{eq:local-input}
\]
がある。ここで$\cC_k$は導来Morita底変更の冪等完備化であり、対象間のHomの
底変更写像はこの同定と両立する。
\end{hypothesis}

$a=n-1$と置く。相対例外列$\OO_0,\ldots,\OO_{a-1}$と、その右直交
\[
 \cK=\langle\OO_0,\ldots,\OO_{a-1}\rangle^\perp\subset\cC
\]
を得ることは以下で証明する。

\begin{theorem}[rotationの自己同値性]\label{thm:rotation}
仮定\ref{hyp:categorical}の下で
\[
 \rho=\LL_{\OO}\circ T:\cK\longrightarrow\cK
\]
はdg自己同値であり、その逆関手は
\[
 \boxed{\ \rho^{-1}=T^{-1}\circ\RR_{\OO}\big|_{\cK}\ }.
 \tag{R}\label{eq:inverse}
\]
評価・余評価のconeから、両合成と恒等関手との自然な同型が得られる。
この主張には$\cK$のsmoothnessを仮定しない。
\end{theorem}

\begin{theorem}[saturation後に従う二つの帰結]\label{thm:summary}
第\ref{sec:projective}節の局所中心三次商$B$とそのproper模型
$\cC_B=\thick\{\OO(j)\}$、$\cK_B=\Perf(\Gamma_R)$を考える。
追加条件
\[
 \Gamma_R\in\Perf(\Gamma_R^{\mathrm{op}}\otL_R\Gamma_R)
 \tag{SM}\label{eq:SM}
\]
を仮定すると、次が成立する。
\begin{enumerate}[label=(\arabic*)]
\item $\cC_B=D^b(\qgr B)$であり、線束生成残余圏は全coherent残余圏と一致する。
\item $R=(R_0)_{\mathfrak m}$、$R_0$が有限型複素整域なら、
$\mathfrak m$を含むある主開集合上へ$(\Perf(\Gamma_R),\rho)$は
smooth properな代数的dg圏と可逆perfect双加群のpairとして降下する。
\end{enumerate}
従って、これら二項を\eqref{eq:SM}とは独立な未解決仮定として重ねて課す必要はない。
\end{theorem}

この二定理と、定理\ref{er:thm:ER}の局所実現定理を合わせたgeneric復元を
定理\ref{thm:generic}に記す。\eqref{eq:SM}自体の証明とは区別する。

\subsection{先行研究との関係と主張の状態}
mutationの基本操作は既知である。ここでの要点は、特殊ファイバーのSerre traceから
局所底上の必要な双対性を作り、逆rotationが残余圏を保つことを実際に証明する点にある。
「properだから任意の関手に随伴がある」とは主張しない。

有界非可換射影導来圏についてはBurke--Stevenson
\cite[Proposition 4.18, Lemmas 6.1--6.2, Theorem 6.4]{BS}を使う。
smooth proper dg代数の有限型底への降下はToën\cite[Théorème 1]{Toen}の定理である。
perfect加群のfiltered colimitに関する性質は同論文F4、
Toën--Vaquié\cite[Lemma 2.10]{TV}による。
これらの既知の定理を本研究の新結果として数えない。
Orlov\cite[Theorem 3.25]{Orlov}のgluing判定は体上の記述なので、
第\ref{sec:fullness}節では必要な方向を一般の可換底上の双加群三角で証明する。

Lin--Zhang\cite[\S2.1, Theorem 2.2, Proposition 2.4]{ERLinZhang}は、幾何的な相対admissible成分について、
同値核の変形障害の消去と代数的correspondenceを扱う。
その入力には相対perfect projectorと底変更の両立がある。
第\ref{er:sec:etale}節では、同じ障害消去原理を実際のperfect双加群のmoduliで定式化し、
非幾何的なtargetとrotation intertwinerを含めて証明する。
同論文をそのまま非幾何的なdg族へ適用したとは扱わない。

本稿のrotation定理と、その後の還元命題には証明を記す。
具体的量子化への適用は第\ref{sec:projective}節の入力の検証に依存する。
全Poisson成分での局所量子化の存在と\eqref{eq:SM}は未解決である。
本稿の組合せとしての新規性は未確定であり、外部査読・証明支援系による形式検証は行っていない。

\section{一つのSerre traceの持上げ}\label{sec:trace}
\begin{lemma}[perfect複体に対する局所Nakayama]\label{lem:nakayama}
$V\in\Perf(R)$が$V\otL_Rk=0$を満たせば$V=0$である。
$V\otL_Rk\simeq k[0]$なら$V\simeq R[0]$である。
\end{lemma}
\begin{proof}
有限自由加群からなる有界複体を取り、微分行列の単元成分を順次消去する。
残った最小複体の微分は全て$\mathfrak m$に値を取り、剰余上では零になる。
剰余が零なら全ての項が零であり、剰余が$k[0]$なら次数零の階数一の項だけが残る。
\end{proof}

\begin{lemma}[例外列と残余生成子]\label{lem:exceptional}
$\OO_0,\ldots,\OO_{a-1}$は相対例外列である。
$\cC=\langle\cK,\OO_0,\ldots,\OO_{a-1}\rangle$を持ち、
\[
 p=\LL_{\OO_0}\circ\cdots\circ\LL_{\OO_{a-1}},\qquad G=p(E)
\]
は$\cK$の生成子である。$\Gamma_R=\REnd(G)$は$R$-perfectである。
\end{lemma}
\begin{proof}
単位写像$R\to\RHom(\OO_i,\OO_i)$のconeと、$j>i$について
$\RHom(\OO_j,\OO_i)$は$R$-perfectで、その剰余は零である。
特殊ファイバーでの消滅は三次超曲面の線束コホモロジーから従う。
補題\ref{lem:nakayama}を適用する。
proper性により、各例外対象への左右のHomはperfectであり、評価・余評価は$\cC$の対象を保つ。
右端から左mutationを繰り返すと指定した右直交への射影$p$を得る。
$E$の生成性から$p(E)$の生成性が従い、proper性から$\Gamma_R$は$R$-perfectである。
\end{proof}

\begin{proposition}[持ち上げたSerre pairing]\label{prop:trace}
ある射
\[
 \operatorname{tr}:\RHom_{\cC}(\OO_a,\OO[n])\longrightarrow R
\]
を選べて、任意の$M\in\cC$について合成とtraceが与える写像
\[
 \beta_M:\RHom_{\cC}(M,\OO[n])
 \longrightarrow\RHom_{\cC}(\OO_a,M)^\vee
 \tag{D}\label{eq:duality}
\]
は擬同型である。従って
\[
 M\in{}^\perp\OO\quad\Longleftrightarrow\quad M\in\OO_a^\perp.
 \tag{O}\label{eq:orthogonals}
\]
\end{proposition}
\begin{proof}
$L=\RHom(\OO_a,\OO[n])$は$R$-perfectである。
$\omega_X=\OO_X(-a)$より、$L\otL_Rk\simeq k[0]$であり、
非零射$L_k\to k$はスカラー倍を除いて$X$のSerre traceである。
補題\ref{lem:nakayama}から$L\simeq R$なので、そのtraceを還元にもつ同型を一つ選ぶ。

合成$\RHom(M,\OO[n])\otL\RHom(\OO_a,M)\to L$とtraceから
\eqref{eq:duality}を作る。両辺が$R$-perfectなのでconeもperfectである。
導来底変更、合成、perfect双対は両立する。剰余上の$\beta_M$は
$S_X(\OO_X(a))=\OO_X[n]$に対するSerre dualityそのものなので擬同型である。
再び補題\ref{lem:nakayama}を使う。
双対はperfect複体上保守的なので\eqref{eq:orthogonals}を得る。
\end{proof}

\begin{remark}
これは対角加群の持上げではない。底上perfectであることが既知のHom複体に対して
一つのtrace射を選び、そのconeにNakayamaを適用したものである。
$\Gamma_R$が包絡代数上perfectであるという主張は含まれない。
\end{remark}

\section{逆rotationの構成}\label{sec:rotation}
$V\in\cC$についてmutationを次のfunctorial triangleで定める：
\begin{align}
 \RHom(\OO,V)\otL_R\OO&\longrightarrow V
       \longrightarrow\LL_{\OO}V,\label{eq:left}\tag{LM}\\
 \RR_{\OO}V&\longrightarrow V
       \longrightarrow\RHom(V,\OO)^\vee\otL_R\OO.
       \label{eq:right}\tag{RM}
\end{align}
proper性により全ての項は$\cC$に属する。

\begin{lemma}[直交圏間のmutation]\label{lem:mutation}
$\LL_{\OO}$と$\RR_{\OO}$は互いに逆な同値
\[
 {}^\perp\OO\ \underset{\RR_{\OO}}{\stackrel{\LL_{\OO}}{\rightleftarrows}}\ \OO^\perp
\]
を与える。
\end{lemma}
\begin{proof}
\eqref{eq:left}に$\RHom(\OO,-)$を適用すると、評価により最初の写像は同型となる。
従って$\LL_{\OO}V\in\OO^\perp$である。
\eqref{eq:right}に$\RHom(-,\OO)$を適用し、perfect二重双対性を使うと
$\RR_{\OO}V\in{}^\perp\OO$となる。
両mutationは$\OO\otL_R\Perf(R)$を零に送る。
$V\in{}^\perp\OO$なら$\RR_{\OO}V\simeq V$なので、
\eqref{eq:left}へ$\RR_{\OO}$を適用して
$\RR_{\OO}\LL_{\OO}V\simeq V$を得る。
同様に$W\in\OO^\perp$について
$\LL_{\OO}\RR_{\OO}W\simeq W$である。
全ての射は評価・余評価から作るため、対象に自然である。
\end{proof}

\begin{proof}[定理\ref{thm:rotation}の証明]
まず$M\in\cK$なら$1\le i<a$に対し
$\RHom(\OO_i,TM)=\RHom(\OO_{i-1},M)=0$である。
$\RHom(\OO_i,\OO)=0$なので、\eqref{eq:left}によるmutation後もこの消滅は保たれる。
$i=0$の消滅は左mutationの定義から従い、$\rho M\in\cK$を得る。
さらに
\[
 \RHom(\OO_a,TM)=\RHom(\OO_{a-1},M)=0
\]
と命題\ref{prop:trace}から$TM\in{}^\perp\OO$である。

次に$N\in\cK$として$Q=\RR_{\OO}N$と置く。
$Q\in{}^\perp\OO$なので命題\ref{prop:trace}から$\RHom(\OO_a,Q)=0$。
$1\le i<a$について\eqref{eq:right}に$\RHom(\OO_i,-)$を適用すると、
$N\in\cK$と$\RHom(\OO_i,\OO)=0$により$\RHom(\OO_i,Q)=0$を得る。
従って$0\le j<a$に対し
\[
 \RHom(\OO_j,T^{-1}Q)=\RHom(\OO_{j+1},Q)=0.
\]
これは$\psi=T^{-1}\RR_{\OO}$が$\cK$を保つことを示す。

補題\ref{lem:mutation}より、$TM\in{}^\perp\OO$と$N\in\OO^\perp$に対して
\[
 \psi\rho M=T^{-1}\RR_{\OO}\LL_{\OO}TM\simeq M,
 \qquad
 \rho\psi N=\LL_{\OO}\RR_{\OO}N\simeq N.
\]
これらはdg自然同型であり、定理が従う。
\end{proof}

\begin{corollary}[底変更]\label{cor:rotation-bc}
Homの導来底変更が成立し、properな模型と$T,\OO$が保たれる底変更について、
上のrotation、逆rotation、両合成の同型は底変更と両立する。
\end{corollary}
\begin{proof}
perfect双対、評価、余評価、有限cone、$T^{\pm1}$はこの底変更と両立する。
逆関手の同型はそれらのfunctorial triangleから作ったので、独立に同型を持ち上げる必要はない。
\end{proof}

\begin{remark}[周期性とは異なる]
本定理は$\rho^3\simeq[2]$を証明しない。
中心三次商のgraded singularity categoryでの$(3)\simeq[2]$と、射影残余圏上の
$\LL_{\OO}(1)$の比較にはwindowの同定が必要である。
特殊ファイバーで同型である二つの関手を、自然変換を構成せずに局所底上で同一視してはならない。
\end{remark}

\section{Serre関手と有限の開条件}\label{sec:serre}
\begin{proposition}[局所Serre関手]\label{prop:serre}
仮定\ref{hyp:categorical}の$\cC$には相対Serre自己同値$S_{\cC}$が存在し、対象として
\[
 S_{\cC}(\OO_i)\simeq\OO_{i-a}[n]\qquad(i\in\mathbb Z)
\]
を満たす。$\cK$は$\cC$で両側admissibleであり、相対Serre自己同値を持つ。
\end{proposition}
\begin{proof}
\eqref{eq:duality}を$T$で移すと
$M\mapsto\RHom(\OO_i,M)^\vee$は$\OO_{i-a}[n]$で表現される。
この表現可能性を持つ第一変数の集合は、有限cone、shift、直和因子で閉じている。
$E$が生成子なので全ての第一変数で表現可能であり、dg YonedaによりSerre準関手を得る。
二重双対性からこれは完全忠実である。その像は生成子$T^{-a}E[n]$を含むので本質的全射である。

有限相対例外列の生成部分は左右随伴を持つ。
Serre関手は左直交を右直交へ移すので、その右直交$\cK$も両側admissibleとなる。
包含$i:\cK\to\cC$の右随伴を$i^!$とすれば、$i^!S_{\cC}i$がそのSerre関手である。
\end{proof}

\begin{remark}
対象$\OO_i$上の同型から、関手として$S_{\cC}=T^{-a}[n]$とは結論しない。
生成対象上の自己Homへの作用を比較する必要があり、Nakayama型のtwistを見落としてはならない。
\end{remark}

\begin{proposition}[代数的proper族の有限trace証明書]\label{prop:open}
有限型noetherianスキーム$S$上に、properなdg圏の族$\cC_S$、$T$、$\OO$、
有限個のtwistからなる生成子$E$が実際に存在するとする。
点$s$で\eqref{eq:local-input}が成立すれば、$s$のZariski開近傍$U$上で
$\LL_{\OO}T$と$T^{-1}\RR_{\OO}$は残余圏の互いに逆な自己同値となる。
$U$は有限個のperfect複体の消滅条件で定められる。
\end{proposition}
\begin{proof}
有限個の単位写像のconeと逆順Homの消滅開集合で縮小し、相対例外列を得る。
$L=\RHom(\OO_a,\OO[n])$は$s$で階数一・振幅零なので、さらに縮小して
次数零の線束となる。$s$の近傍で自明化してtraceを選ぶ。
有限のperfect複体
\[
 W=\Cone\!\left(\RHom(E,\OO[n])
              \xrightarrow{\beta_E}\RHom(\OO_a,E)^\vee\right)
 \tag{FC}\label{eq:finite-certificate}
\]
は$s$で零である。そのコホモロジーのsupportを除くと$W=0$になる。
$\beta_M$が擬同型となる$M$の全体はthickなので、$E$の生成性から全ての$M$で成り立つ。
あとは第\ref{sec:rotation}節の証明をそのまま適用できる。
\end{proof}

\begin{remark}
\eqref{eq:finite-certificate}は有限の圏論的証明書の\emph{形式}であり、
具体的$\Gamma_R$の行列を書き下して数値計算したという意味ではない。
この命題はproper族が有限型底上に既にある場合の主張であり、その族自身の降下を代替しない。
\end{remark}

\section{局所中心三次商からの入力}\label{sec:projective}
ここでは本稿で使う既存局所模型を、具体的な代数的仮定まで遡って記す。
本節は親リポジトリ\cite{Repo}の\texttt{local\_projective\_model.tex}と同じ範囲の結果であり、
本稿の新しいrotation論証と区別する。必要な推論を以下に再構成する。

\begin{hypothesis}[次数別flatな局所量子化]\label{hyp:algebra}
$R$は有限型複素代数の局所化である正則局所環、$d=\dim R$、$N=n+2\ge5$とする。
$A=\bigoplus_{j\ge0}A_j$は中央$R$-代数、$A_0=R$で、各$A_j$は有限自由であり、
$A_k\simeq k[z_1,\ldots,z_N]$である。
$\Omega\in Z(A)_3$の剰余$F$は滑らかな三次超曲面を定める。
$B=A/(\Omega)$、$a=N-3$とする。
\end{hypothesis}

\begin{proposition}[局所proper模型の入力検証]\label{prop:local-model}
仮定\ref{hyp:algebra}の下で$A,B$は左右noetherianで、$\Omega$は中心正則元である。
$\cC_B=\thick_R\{\OO(j):j\in\mathbb Z\}\subset D^b(\qgr B)$は仮定\ref{hyp:categorical}を満たす。
生成子は$E=\bigoplus_{i=0}^{N-1}\OO(i)$でよい。
また全ての$M\in D^b(\qgr B)$について$\RHom(E,M)$は$R$-perfectであり、$E$はこの大きい圏で弱生成子である。
\end{proposition}

\begin{proof}
\emph{(i) noetherian性と自由分解。}
一般に$D_j$が有限$R$-加群で$D/\mathfrak mD$が右noetherianなら、
$\operatorname{gr}_{\mathfrak m}D$は$(D/\mathfrak mD)[t_1,\ldots,t_e]$の商なので右noetherianである。
斉次右イデアル$I$の随伴次数付きイデアルの有限個の二重斉次生成元を持ち上げ、
それらの生成するイデアルを$J$とする。最低$\mathfrak m$次数を順に消すと
$I_j\subset J_j+\mathfrak m^qD_j$を全ての$q$で得る。
有限$R$-加群$D_j/J_j$にKrull交叉を適用して$I_j=J_j$となる。
任意の右イデアルは内部次数の最高項を用いて同様に有限生成となる。左側も同じである。

$D,M$が次数別有限自由で、剰余上の$M_k$が有限階数自由加群による有限graded分解を持つとする。
第零項の生成元の像を持ち上げると、次数別Nakayamaから全射となる。
その核は各次数で分裂直和因子なので有限自由であり、剰余上の核も正しく得られる。
これを反復すると同じshiftと階数の分解を持ち上げられる。最後の核は次数別Nakayamaで零である。
この操作を多項式環のKoszul分解と対角Koszul分解に適用し、
$R$の$A$-分解および$A$の$A^e$-分解を得る。項はそれぞれ
$A(-j)^{\binom Nj}$と$A^e(-j)^{\binom Nj}$、$0\le j\le N$である。
全分解は次数別$R$-splitなので任意の通常の底変更で完全性を保つ。
対角分解を加群とテンソルすれば$\operatorname{gldim}A\le N+d$となる。
次数一生成と二次関係も、剰余上からの同じ持上げで従う。

\emph{(ii) Gorenstein性と中心商。}
Koszul分解を$A$で双対化すると、各内部次数は有界有限自由$R$-複体である。
剰余上の完全性と補題\ref{lem:nakayama}から、内部次数$-N$、cohomological degree $N$以外は零となる。
この唯一の箇所では$R$が残る。従って左右で
\[
 \operatorname{R\underline{Hom}}_A(R,A)\simeq R(N)[-N].
\]
$\Omega:A_{j-3}\to A_j$の剰余は単射で、列数大の小行列式に単元がある。
従って各次数で分裂単射であり、$\Omega$は正則、$B_j$は有限自由となる。
長さ一の$A$-分解からchange-of-rings
\[
 \operatorname{R\underline{Hom}}_A(M,A)
 \simeq\operatorname{R\underline{Hom}}_B(M,B)(3)[-1]
\]
を得る。$M=R$を代入すると
$\operatorname{R\underline{Hom}}_B(R,B)=R(a)[-(N-1)]$である。
$A$の有限大域次元から$B$の左右injective dimensionは高々$N+d-1$となる。
$B$は商なのでnoetherianである。

\emph{(iii) 線束Homとその底変更。}
Burke--Stevenson\cite[Lemmas 6.1--6.2]{BS}により$B$は$\chi$と有限cohomological dimensionを満たす。
$M_\ell=B/B_{\ge\ell}$を内部次数で濾過すると、その商は$B_j\otimes_RR(-j)$である。
上のGorenstein型から、これらの$B$へのExtは次数$N-1$、内部次数$-a-j$にのみ存在する。
各固定内部次数$t$では$j=-t-a$という唯一の層が寄与し、$\ell>j$で安定する。
従って局所コホモロジーは
\[
 H^q(R\tau B)_t=
 \begin{cases}B_{-t-a}^\vee&q=N-1,\\0&q\ne N-1.\end{cases}
\]
三角$R\tau B\to B\to R\omega\pi B$から
\[
 \operatorname{Ext}^p(\OO(i),\OO(j))=
 \begin{cases}
 B_{j-i}&p=0,\\ B_{i-j-a}^\vee&p=n,\\0&\text{その他}
 \end{cases}
 \tag{H}\label{eq:line-hom}
\]
となる。$a>0$なので二群が同じ差で同時に非零となることはない。

底変更はコホモロジー階数の一致だけでなく、自然な導来Hom写像について確認する。
$M_\ell$の有限階数graded自由分解を取り、各内部次数で$R$-splitにする。
$\underline{\operatorname{Hom}}_B(P_\bullet,B)_t$は下に有界な有限自由複体で、
そのコホモロジーは一つの次数の有限自由加群である。
$R$は有限大域次元なので、完全な射影複体のcycleにdimension shiftingを適用すると
完全部分は分裂する。従って通常のテンソルが導来テンソルを計算する。
項ごとのHom底変更と、各$t$で安定するfiltered colimitを使えば、線束Homの自然な底変更が従う。
この議論は$R'\!$が正則noetherian・有限大域次元で$A\otimes_RR'$がnoetherianである
通常の底変更に適用され、特に剰余体と分数体に適用される。

\emph{(iv) 有限生成と特殊ファイバー。}
Burke--Stevenson\cite[Theorem 6.4]{BS}を$A$に適用すると、$\Dsg(\gr A)=0$より
任意の連続した$N$個のtwistが$D^b(\qgr A)$を生成する。
$A\to B$は有限Tor振幅を持ち、導来制限は$\OO_A(j)$を$\OO_B(j)$へ送る。
よって全twistは$E$のthick閉包に入る。
\eqref{eq:line-hom}から$\REnd(E)$は$R$-perfectで、有限cone・直和因子によって$\cC_B$全体がproperとなる。
自然なHom底変更は積と両立し、有限cone・冪等完備化へ延長できる。
特殊ファイバーで全twistは$\Perf(X)$を生成するため、\eqref{eq:local-input}を得る。

\emph{(v) 全coherent圏へのHom。}
\cite[Proposition 4.18]{BS}の有界導来sectionは$D^b(\qgr B)$を$D^b(\gr_{\ge i}B)$へ送る。
各内部次数の成分は有限$R$-加群である。従って固定した$j$について
$\RHom(\OO(j),M)$は有界かつ有限生成コホモロジーを持つ。
$R$の正則性よりこれは$R$-perfectである。
$\RHom(E,M)=0$なら全twistからのHomが零であり、導来sectionの各次数が零になる。
同命題の完全忠実性から$M=0$となる。
\end{proof}

\begin{corollary}[局所量子化のrotation]\label{cor:local-rotation}
仮定\ref{hyp:algebra}の残余部分圏$\cK_B$上で
$\rho=\LL_{\OO}(1)$は自己同値であり、$\rho^{-1}=(-1)\RR_{\OO}$である。
残余smoothnessは仮定しない。
\end{corollary}

\section{smoothnessと全coherent残余圏}\label{sec:fullness}
\begin{lemma}[smooth代数上のpseudo-perfect加群]\label{lem:pseudoperfect}
$H$は可換環$R$上smoothなdg代数とする。
$H$-加群$V$の基礎$R$複体がperfectなら、$V$は$H$-perfectである。
\end{lemma}
\begin{proof}
対角$H$は$H^e$の有限cone・shift・直和因子から作られる。
これに$V\otL_H-$を作用させると、自由双加群の像は$V\otL_RH$である。
$V$が$R$-perfectなのでこれは$H$-perfectであり、その有限操作で$V$が得られる。
これは\cite[Lemma 2.8(2)]{TV}、\cite[F6]{Toen}の性質の直接の証明である。
\end{proof}

\begin{lemma}[弱生成子による全体の回収]\label{lem:fullness}
$\cD$を冪等完備dg圏、$E\in\cD$を弱生成子とし、$H=\REnd(E)$と置く。
全ての$M\in\cD$について$\RHom(E,M)$が$R$-perfectで、$H$が$R$上smoothなら
$\cD=\thick_R(E)$である。
\end{lemma}
\begin{proof}
$V=\RHom(E,M)$は補題\ref{lem:pseudoperfect}により$H$-perfectである。
$\thick(E)\simeq\Perf(H)$から、それに対応する$N\in\thick(E)$を選ぶ。
任意の$N\in\thick(E)$について、自然な写像
\[
 \RHom_{\cD}(N,M)\longrightarrow
 \RHom_H(\RHom(E,N),\RHom(E,M))
\]
は擬同型である。これは$N=E$のdg Yonedaから有限操作で延長すればよい。
従って選んだ加群同型を射$N\to M$へ戻せる。
そのconeは$E$からの全Homが零なので、弱生成性により零となる。
\end{proof}

\begin{proposition}[相対gluingのsmoothness]\label{prop:gluing}
$R$-properで有限生成なdg圏$\cC$が、両側admissibleな半直交分解
$\cC=\langle\cK,\OO_0,\ldots,\OO_{a-1}\rangle$を持つとする。
$\OO_i$が相対例外対象なら、$\cC$が$R$上smoothであることと$\cK$が$R$上smoothであることは同値である。
\end{proposition}
\begin{proof}
まず二成分のgluingを考える。生成子を取り、零である方向のHom複体を零に置き換えると、
次の三角dg代数にMorita同値な模型を得る：
\[
 H=\begin{pmatrix}A&M\\0&D\end{pmatrix},\qquad e+f=1.
\]
$A,D$がsmoothで$M$が$R$-perfectなら、$A\otL_RD^{\mathrm{op}}$もsmoothなので
補題\ref{lem:pseudoperfect}から$M$は$A$-$D$双加群としてperfectである。
次の対角の三角を用いる：
\[
 He\otL_A M\otL_D fH\longrightarrow
 (He\otL_A eH)\oplus(Hf\otL_D fH)
 \longrightarrow H.
 \tag{G}\label{eq:gluing-diagonal}
\]
最初の射は$m$を左右どちら側に掛けるかの差、二番目は乗法である。
homotopy-flatな模型を取り各$e,f$成分を見れば、対角成分は恒等射、
交差成分は$M\to M\oplus M\to M$のsplit完全列であり、この三角が確認できる。

$A$の対角の有限perfect表示を両側から誘導すれば$He\otL_A eH$は$H^e$-perfectとなる。
$D$の項も同様である。$M$のperfect性から第一項も$H^e$-perfectなので、$H$の対角はperfectとなる。
各例外成分$\Perf(R)$を一つずつ加える帰納法により、$\cK$のsmoothnessから$\cC$のsmoothnessを得る。
この議論は体を仮定せず、Orlov\cite[Theorem 3.25]{Orlov}の必要な方向を相対化したものである。

逆に$\cC$がsmoothなら、その有限perfect対角を$\cK$の両変数へ制限する。
左右随伴により、制限された表現可能加群は再び表現可能である。
したがって制限はperfect性を保ち、$\cK$の対角もperfectとなる。
\end{proof}

\begin{proof}[定理\ref{thm:summary}(1)の証明]
命題\ref{prop:serre}により$\cK_B$は両側admissibleである。
\eqref{eq:SM}と命題\ref{prop:gluing}から$\cC_B$がsmoothとなる。
命題\ref{prop:local-model}の最後の主張と補題\ref{lem:fullness}により
$\cC_B=D^b(\qgr B)$となる。
同じ例外列の右直交を取れば、線束生成残余圏と全coherent残余圏が一致する。
既存の有界Orlov分解との同定も、この右直交を用いるwindowで得られる。
ただし任意の導来底変更に関するOrlov核のcoherenceまで主張しているのではない。
\end{proof}

\section{smooth proper pairの有限型底への降下}\label{sec:spread}
\begin{proposition}[局所pairのspread]\label{prop:spread}
$R=(R_0)_{\mathfrak m}$を有限型複素整域の局所化とする。
$\Gamma_R$が$R$上smooth properで、$\rho$が$\Perf(\Gamma_R)$の$R$-linear dg自己同値なら、
ある$f\in R_0\setminus\mathfrak m$と$U=\operatorname{Spec}R_0[f^{-1}]$上の
smooth proper dg代数$\Gamma_U$、可逆perfect双加群$P_U$が存在し、
局所化後に$(\Perf(\Gamma_R),\rho)$を与える。
\end{proposition}
\begin{proof}
$R=\varinjlim_{f\notin\mathfrak m}R_0[f^{-1}]$である。
Toën\cite[Théorème 1]{Toen}により$\Gamma_R$は、ある段階のsmooth proper dg代数へ降下する。
$\rho$とその逆関手を表すMorita双加群を$P,Q$とする。
その基礎$R$複体は、生成子とその像のHomで計算できるのでperfectである。
$\Gamma_R^e$もsmoothであるから、補題\ref{lem:pseudoperfect}により
$P,Q\in\Perf(\Gamma_R^e)$である。

perfect加群のfiltered colimit性\cite[F4]{Toen}\cite[Lemma 2.10]{TV}により
$P,Q$と逆同型
\[
 P\otL_{\Gamma_R}Q\simeq\Gamma_R,
 \qquad Q\otL_{\Gamma_R}P\simeq\Gamma_R
\]
を有限の段階へ降下できる。射、逆射、および有限個の合成が恒等となる関係も、
同じfiltered colimit性を使ってさらに一段階先で成立させられる。
従って降下した$P_U$は可逆なperfect双加群である。
元の$\rho$との同定は$R$上で保たれる。
\end{proof}

\begin{proof}[定理\ref{thm:summary}(2)の証明]
系\ref{cor:local-rotation}により$\rho$は実在するdg自己同値である。
\eqref{eq:SM}と既存のproper性の下で命題\ref{prop:spread}を適用する。
\end{proof}

\begin{remark}
この証明は、$R$-perfectなdg代数を有限階数のstrict dg代数と同一視しない。
個々の高次$A_\infty$積の分母を無限に列挙して一つの開集合を得ようとする議論も不要である。
smoothnessが与えるhomotopical finite presentationが有限型降下を可能にする。
\end{remark}

\section{同値双加群の障害消去とétale局所実現}\label{er:sec:etale}
本節では、可換模型の存在に使う局所実現原理を証明する。
入力は実在するsmooth properな代数的dg族であり、proper性だけの族ではない。
圏の全変形をHochschild複体から実現するという主張は用いず、
実在するperfect双加群のmoduliと、その完全な障害理論を用いる。
ここでのpairは自己同値一つを指定したもので、$\rho^3\simeq[2]$という
選択された同型やそのcoherenceを追加データとして含めない。

\subsection{perfect加群の小拡大障害}
小拡大は剰余体$k=\CC$を持つ局所Artin代数の全射
$A'\twoheadrightarrow A$で、核$I$が$\mathfrak m_{A'}I=0$を満たすものとする。
$D'$は$A'$上smooth properなdg代数とし、$D=D'\otL_{A'}A$とする。

\begin{lemma}[加群の完全な障害理論と特徴作用]\label{er:lem:obstruction}
$P\in\Perf(D)$の持上げには自然な障害
\[
 o(P;D')\in\operatorname{Ext}^2_{D_k}(P_k,P_k)\otimes_k I
 \tag{O}\label{er:eq:obstruction}
\]
があり、零であることと$\Perf(D')$への持上げの存在は同値である。
障害はMorita関手と対象の同型に関して自然である。
$D'=C'^{\rm op}\otL_{A'}B'$で$P$が$C$から$B$への同値を表すとき、
$B'$を固定し、$C'$のliftをHochschild類$\xi\in\mathrm{HH}^2(C_k)\otimes I$だけ変えると、
障害の変化は、符号を一貫して選べば
\[
 o_{\rm new}-o_{\rm old}=-\chi_P(\xi),\qquad
 \chi_P:\mathrm{HH}^2(C_k)\xrightarrow{\sim}
       \operatorname{Ext}^2_{C_k^{\rm op}\otimes B_k}(P_k,P_k).
 \tag{V}\label{er:eq:variation}
\]
\end{lemma}
\begin{proof}
Toën--Vaquié\cite[Theorem 3.6, Corollaries 3.15, 3.17]{TV}を
可換底$A'$と$D'$に適用する。smoothnessによりpseudo-perfect加群と
perfect加群は一致し、その実際の対象moduliの相対接複体は$\RHom_D(P,P)[1]$である。
導来対象moduliの平方零拡大に関する障害理論は、その接複体の次数一に完全な障害を持つ。
従って障害は$\operatorname{Ext}^2_D(P,P\otL_A I)$に入り、
導来随伴と$\mathfrak m_A I=0$によって\eqref{er:eq:obstruction}に一致する。
これは単なる接空間の計算を代数化の代用にしているのではない。
実在する対象の導来moduliが無限小凝集的であることを使っており、
障害のnull-homotopyはそのmoduliの$A'$-点、すなわち実際のperfect加群を与える。

変化式も説明する。cofibrantな代数模型と加群のbar模型を選び、
構造を次数一のbar余微分$b$で表す。任意のlift$\widetilde b$の曲率
$\widetilde b^{\,2}$は$I$に値を取り、$[b,\widetilde b^{\,2}]=0$である。
持上げを$I$値の$h$だけ変えると、$I^2=0$により曲率は$[b,h]$だけ変わる。
この曲率類が上のmoduliの障害である。
代数側を$\xi$だけ変えた際の曲率の変化は、$\xi$を加群の構造写像に挿入したcochainであり、
それが特徴作用$\chi_P(\xi)$である。符号はsourceの反対代数とtorsorの規約に合わせて
\eqref{er:eq:variation}と取る。必要なのは同じ規約を両式で使うことである。
この記述はLowen\cite[Theorem 4.8]{ERLowen}の特徴作用による障害の記述を、
bar模型でdg加群と任意の小拡大に書いたものである。
$P_k$の逆Morita双加群をテンソルすれば、その自己Homと対角の自己Homが同定されるので、
$\chi_P$は同型である。曲率の構成、またはmoduliの障害の自然性から、最後の自然性も従う。
\end{proof}

\begin{lemma}[還元後の同値とintertwiner]\label{er:lem:lift}
$C',B'$は$A'$上smooth properとする。
$P'\in\Perf(C'^{\rm op}\otL B')$の還元がMorita同値なら$P'$自身も同値である。
さらに$Q'$もperfectで、$\operatorname{Ext}^1(P_k,Q_k)=0$なら、
任意の射$P'_A\to Q'_A$は$A'$上へ持ち上がる。元の射が同型ならliftも同型となる。
\end{lemma}
\begin{proof}
perfect双加群は両側に随伴を持ち、smooth proper性により随伴双加群もperfectである。
unitとcounitのconeは底上perfectで、還元後に零であるから導来Nakayamaで零となる。
射の持上げには三角$Q'\otL I\to Q'\to Q'_A$を$\RHom(P',-)$で評価する。
射の障害は$\operatorname{Ext}^1(P_k,Q_k)\otimes I$に入り、仮定で零である。
持ち上げた射のconeに同じNakayamaを適用すれば同型性も従う。
\end{proof}

\subsection{rotationを保つ障害消去}
$H$は複素数体上smoothな有限型スキーム、$h\in H(\CC)$とする。
$\mathcal A_H$はsmooth proper dg族、$r_H$は可逆perfect双加群で表された自己同値とする。
$\mathcal A=\mathcal A_h$、$r=r_h$と置く。局所的なdg代数模型を取って議論する。

\begin{theorem}[自己同値付き局所実現の判定]\label{er:thm:criterion}
次を仮定する。
\[
 \mathrm{HH}^{<0}(\mathcal A)=0,\qquad
 \mathrm{HH}^{1}(\mathcal A)=0,
 \qquad
 T_hH\longrightarrow\mathrm{HH}^2(\mathcal A)^r
 \text{ が全射}.
 \tag{KS}\label{er:eq:KS}
\]
最後の写像は$\mathcal A_H$のKodaira--Spencer写像である。
$S$は有限型複素スキーム、$(\mathcal B_S,s_S)$はsmooth proper dg族と
可逆perfect双加群のpairとする。$s\in S(\CC)$に指定した同値
$(\mathcal A,r)\simeq(\mathcal B_s,s_s)$があるなら、
pointed étale近傍$V\to(S,s)$、$f:V\to H$、$f(v)=h$と
指定同値を延長するpairの同値
\[
 (f^*\mathcal A_H,f^*r_H)\simeq(\mathcal B_V,s_V)
\]
が存在する。
\end{theorem}
\begin{proof}
まず小拡大$A'\twoheadrightarrow A$上の持上げを証明する。
$\operatorname{Spec}A'\to S$を固定し、$A$上の点$h_A$、同値双加群$P_A$、
intertwiner $\alpha_A:s_AP_A\xrightarrow{\sim}P_Ar_A$が与えられているとする。
$H$はsmoothなので$h_A$の任意のlift$h'_0$を取れる。
両側のliftされた自己同値から、双加群圏の自己同値
\[
 \mathcal F'(P)=s'\circ P\circ(r'_0)^{-1}
\]
を得る。$\alpha_A$は$\mathcal F_A(P_A)\simeq P_A$を与える。
補題\ref{er:lem:obstruction}の自然性より、$P_A$の持上げ障害$o$は
$\mathcal F_k$の作用で不変である。
Morita同定$\operatorname{Ext}^2(P_k,P_k)=\mathrm{HH}^2(\mathcal A)$の下で、
この作用は$r$による共役作用（規約によりその逆）である。
従って
\[
 o\in\mathrm{HH}^2(\mathcal A)^r\otimes I.
 \tag{I}\label{er:eq:invariant}
\]
この段階では$\alpha_A$を$A'$上へ持ち上げたとは仮定していない。

\eqref{er:eq:KS}と\eqref{er:eq:variation}により、$h'_0$を
$T_hH\otimes I$の元だけ修正して$o$を消せる。
修正後のsourceと固定したtargetの間にperfect双加群$P'$が存在し、
補題\ref{er:lem:lift}によって同値となる。
$\alpha_A$を持ち上げる障害は
\[
 \operatorname{Ext}^1(s_kP_k,P_kr_k)\otimes I
       \simeq \mathrm{HH}^1(\mathcal A)\otimes I=0
\]
なので、intertwinerも同型として持ち上がる。
自己同値一つのintertwiner以外に群関係は課していないので、ここで必要なデータは全て揃う。
任意のnilpotent拡大は小拡大に分解できる。

次に代数化する。$H\times S$上で、perfectな
$\mathcal A_H^{\rm op}\otimes\mathcal B_S$-加群$P$と
同型$s_SP\simeq Pr_H$のmoduliを考え、$P$が同値である部分を$\mathcal E$とする。
Toën--Vaquié\cite[Theorem 3.6, Corollaries 3.15, 3.19]{TV}により、
perfect対象とその射の導来moduliは局所幾何的かつ局所有限表示である。
unit、counitおよびintertwinerのconeは底上perfectだから、その消滅は開条件である。
従って$\mathcal E$も局所幾何的かつ局所有限表示である。
同値点での自己Extの負次数は$\mathrm{HH}^{<0}(\mathcal A)=0$であり、
その近傍でこの消滅を保てる。従って古典的切断はその近傍でArtin 1-stackとなる。
ここでは局所dg代数模型を用いてTVの定理を$H\times S$上の可換係数環に適用している。

上の小拡大持上げは、指定点$e=(h,s,P_k,\alpha_k)$で
$\mathcal E\to S$が形式的にsmoothであることを示す。
局所有限表示のArtin stackの無限小判定により、この射は$e$でsmoothである。
$e$を覆うsmooth atlas $U\to\mathcal E$を取り、点の近傍へ縮めると$U\to S$もsmoothである。
滑らかな射のpointed étale局所section\cite[Tags 02HX, 055U]{StacksER}を使い、
$V\to S$と$V\to U\to\mathcal E$を得る。
この普遍対象が所要の族、双加群およびintertwinerを与える。
\end{proof}

\begin{remark}[使わない仮定と先行研究]\label{er:rem:etale-scope}
この証明は$\mathrm{HH}^3=0$、Hochschild cochainのformality、
抽象的なHochschild変形の全てがactual Morita変形になること、または$\rho^3=[2]$の指定を使わない。
Lin--Zhang\cite[Theorem 2.2, Proposition 2.4]{ERLinZhang}の同値核の障害消去と同じ原理を用いるが、
ここではtargetを非幾何的なsmooth proper dg族として扱い、
自己同値が障害を不変部分へ制限することと、intertwinerの持上げを別々に証明した。
同論文の幾何的projectorの仮定を満たさないtargetへその定理を直接適用していない。
この一般定式化の新規性は未確認である。
\end{remark}

\subsection{滑らかなcubicへの適用}
\begin{lemma}[低次Hochschild群と幾何的接写像]\label{er:lem:cubic}
$X=V(F)\subset\mathbb P^{N-1}$を滑らかなcubic、$N=n+2\ge5$とし、
$J(F)=\CC[z_1,\ldots,z_N]/(\partial_1F,\ldots,\partial_NF)$とする。
\[
 \mathrm{HH}^{<0}(\Ku(X))=0,\quad
 \mathrm{HH}^0(\Ku(X))=\CC,\quad \mathrm{HH}^1(\Ku(X))=0,
\]
\[
 \mathrm{HH}^2(\Ku(X))=
 \begin{cases}J(F)_3&n\ne4,\\J(F)_3\oplus\CC^2&n=4,\end{cases}
 \qquad
 H^1(X,T_X)\xrightarrow{\sim}\mathrm{HH}^2(\Ku(X))^{\rho_X}=J(F)_3.
 \tag{H2}\label{er:eq:HH}
\]
$\dim J(F)_3=\binom N3$である。
\end{lemma}
\begin{proof}
BFK\cite[Theorem 5.39]{BFKER}のequivariant matrix factorizationの公式を用いる。
$G=\mathbb G_m$、各座標の重みは1、potentialのcharacterは3である。
恒等sectorのKoszul cohomologyは、偏微分が正則列であるためJacobian商だけで、
cohomological degree $2r$に$J(F)_{3r}$を与える。
非恒等$g\in\mu_3$の固定空間は原点、conormalの次元は$N$、determinantの重みは$-N$である。
次数$q$の重み零条件は$(3q-N)/2=0$であり、従って$3\mid N$かつ$q=N/3$のときだけ
各sectorから一次元が現れる。これで負次数、次数0、1、2の群が得られる。

幾何的Kodaira--Spencer写像の同定は、単なる次元一致ではない。
$n\ge4$では$3\le(n+2)/2$なので、Huybrechts--Rennemo
\cite[Lemma 2.19, \S2.5, \S3.3]{HRER}のAtiyah--Kodaira--Spencerと
三次周期性の比較を適用でき、$F+\epsilon G$の像は恒等sectorの$[G]\in J(F)_3$である。
$n=3$ではJacovskis--Lin--Liu--Zhang\cite[Theorem 4.4(2)]{JLLZER}の
幾何的写像の同型を用いる。normal sequenceとEuler sequenceから$H^1(T_X)=J(F)_3$である。
幾何的族にはrotationがあるのでこの像は不変である。
$n=4$で余分な二方向が不変でないことは、Huybrechts\cite[Proposition 3.12]{HuyER}による。
rotationはsymplecticで、Mukai格子の$A_2$上で三つの根を巡回する。
不変なCalabi--Yau類による$\mathrm{HH}^2\simeq\mathrm{HH}_0$の同定の下で、
その二方向の固有値は$\zeta,\zeta^2$である。よって不変部分は正確に$J(F)_3$となる。
最後に偏微分が二次正則列であることから$H_{J(F)}(t)=(1+t)^N$を得る。
\end{proof}

\begin{theorem}[cubicのrotation付きétale局所実現]\label{er:thm:ER}
$S$は有限型複素スキーム、$\mathcal K_S$は実在するsmooth properな代数的dg族で、
$\rho_S$は可逆perfect双加群により表される自己同値とする。
$s\in S(\CC)$で、$n\ge3$の滑らかなcubic $X$に対する指定同値
$(\mathcal K_s,\rho_s)\simeq(\Ku(X),\rho_X)$があるなら、
pointed étale近傍$V\to(S,s)$、滑らかなcubic族$Y/V$、および指定同値を延長する同値
\[
 (\mathcal K_V,\rho_V)\simeq(\Ku(Y/V),\rho_{Y/V})
 \tag{ER}\label{er:eq:ER}
\]
が存在する。
\end{theorem}
\begin{proof}
滑らかな三次形式の射影パラメータ空間の$[F]$を含むアフィン開を$H$とする。
普遍cubic $\mathcal X/H$のrelative perfect圏はsmooth properで、
$\OO,\ldots,\OO(n-2)$のrelative exceptional列の射影により
smooth properな相対Kuznetsov成分を得る。
射影とrotationは有限個の評価三角で構成され、底変更に可換である。
普遍cubic上の有限個のtwistから生成子を選べるので、局所的なdg代数模型も存在する。
$T_{[F]}H\to H^1(T_X)=J(F)_3$は全射である。
補題\ref{er:lem:cubic}と定理\ref{er:thm:criterion}を適用する。
\end{proof}

\begin{corollary}[rotationを最初に指定しない存在]\label{er:cor:unpolarized}
$n\ne4$なら、前定理で$\rho_S$を入力として指定せずとも、
$\mathcal K_s\simeq\Ku(X)$からétale局所的な可換cubic模型が得られる。
これはrotationなしの大域的な一意性の主張ではない。
\end{corollary}
\begin{proof}
補題\ref{er:lem:cubic}では幾何的接写像が$\mathrm{HH}^2$全体に全射である。
定理\ref{er:thm:criterion}の証明から不変性とintertwinerの段階を省けばよい。
\end{proof}

\section{Hochschildの奇数次障害とrotation不変部分}\label{ho:sec:parity}
可換復元における対角smoothnessを、形式変形の障害と混同しないために、
滑らかな三次超曲面のKuznetsov成分の低次Hochschild群を全次元で整理する。
以下は既知のequivariant matrix factorizationの計算からの帰結であり、
残余dg圏の相対smoothnessに関する新たな一般定理ではない。

\begin{proposition}[奇数次Hochschild群とrotation固有値]\label{ho:prop:sectors}
$X=V(F)\subset\mathbb P^{N-1}_\CC$を滑らかな三次超曲面、$N=n+2\ge5$とし、
$J(F)=\CC[z_1,\dots,z_N]/(\partial F)$とする。
$\operatorname{Ku}(X)$の標準rotationを$\rho$とする。
$\zeta$を原始三乗根と書けば、次の複素ベクトル空間の分解が成り立つ：
\[
 \operatorname{HH}^{q}(\operatorname{Ku}(X))\cong
 \begin{cases}
 J(F)_{3q/2},&q\text{ 偶数},\\ 0,&q\text{ 奇数}
 \end{cases}
 \quad\oplus\quad
 \begin{cases}
 \CC_{\zeta}\oplus\CC_{\zeta^2},&3\mid N,\ q=N/3,\\
 0,&\text{それ以外}.
 \end{cases}
 \tag{HS}\label{ho:eq:sectors}
\]
ここで$J(F)_d=0$（$d<0$）とし、右側の二つの一次元空間は
$\rho$の共役作用の非自明な固有空間である。
$J(F)$から来る項では$\rho$は恒等的に作用する。
特に
\[
 \operatorname{HH}^3(\operatorname{Ku}(X))
 \cong\begin{cases}\CC^2&n=7,\\0&n\ne7,\end{cases}
 \qquad
 \boxed{\operatorname{HH}^3(\operatorname{Ku}(X))^\rho=0\quad(n\ge3).}
 \tag{V3}\label{ho:eq:HH3}
\]
さらに$\operatorname{HH}^2(\operatorname{Ku}(X))^\rho\cong J(F)_3$で、
$\dim_\CC J(F)_3=\binom N3$である。
\end{proposition}
\begin{proof}
Orlovのgraded matrix factorization表示と
Ballard--Favero--Katzarkov\cite[Theorem 5.39, pp.~82--85]{BFKER}の
固定sector公式を用いる。
$G=\mathbb G_m$は全生成変数の重みが$1$となる作用で、
potentialのcharacterは三次である。$K_\chi=\mu_3$とする。
$g=1$において$V_g=\mathbb A^N$、$W_g=0$である。
偏微分$(\partial F)$は原点にのみ共通零点を持つ長さ$N$の正則列であり、
Koszul cohomologyは零次の$J(F)$に集中する。
重み零の条件とcohomological gradingから、恒等sectorは偶数次数$q=2r$に
$J(F)_{3r}$を与え、奇数次数には寄与しない。

$g\in\mu_3\setminus\{1\}$なら固定空間は原点で、$\dim W_g=N$、
$J(F_g)=\CC$である。
Theorem 5.39のtwistとdeterminant characterを代入すると、
$u=(q-N)/2$に対する重み零条件は$3u+N=0$となる。
従って$q=N/3$でのみ各非自明sectorが一次元を与える。
各sectorは標準的な文字twist $(1)$ による共役作用で
$g$の文字（規約によって逆文字）を受ける。
これは同論文のkernelの固定graph分解と
\cite[Theorem 6.5]{BFKER}のcharacter作用の計算と同じものである。
Orlov同値の下で次数shift $(1)$ は$\rho$に対応するため、
二つのsectorの固有値は順序を除いて$\zeta,\zeta^2$となる。
$J(F)$から来る恒等sectorではこの作用は自明である。
これで\eqref{ho:eq:sectors}を得る。

$q=3$は奇数なので、非零になり得るのは$N/3=3$、すなわち$N=9$
（$n=7$）だけである。しかもその二方向には$\rho$-不変ベクトルがない。
$q=2$の不変部分は常に$J(F)_3$である。
滑らかさから偏微分は二次正則列なので、
$H_{J(F)}(t)=(1+t)^N$となり$\dim J(F)_3=\binom N3$を得る。
\end{proof}

\begin{corollary}[非楕円cubic fivefoldの変形障害の所在]\label{ho:cor:fivefold}
本研究の節点型・尖点型の二つの滑らかな可換特殊ファイバー
$X^5$について、共に
\[
 (\dim\operatorname{HH}^0,\dim\operatorname{HH}^1,
 \dim\operatorname{HH}^2,\dim\operatorname{HH}^3,
 \dim\operatorname{HH}^4)=(1,0,35,0,7).
 \tag{F5}\label{ho:eq:fivefold}
\]
従って滑らかな可換特殊ファイバーのdg圏自身の通常のMorita変形に対する
\emph{一次の}Kuranishi障害空間$\operatorname{HH}^3$は零である。
ただし、量子化から構成した$\Gamma_R$の
$\Gamma_R^e$上の対角perfect性を、これだけからは結論できない。
\end{corollary}
\begin{proof}
$N=7$を命題\ref{ho:prop:sectors}に代入する。
$J(F)_0,J(F)_3,J(F)_6$の次元はそれぞれ$1,35,7$で、
それ以外の表示次数$q=1,3$は零となる。
Morita変形の一次障害が$\operatorname{HH}^3$に入ることは
標準的なHochschild変形理論の適用である。
\end{proof}

\begin{remark}[残余対角との論理的境界]\label{ho:rem:barrier}
$\operatorname{HH}^3(\Ku(X))^\rho=0$は重要な形式的情報であるが、
この不変部分の消滅だけで任意のrotation付きdg圏変形の厳密な
$\mathbb Z/3$-equivariant obstruction theoryを構成したとはしない。
また、Artin切断上の変形が可換的に説明できたとしても、
有限型係数環上の\emph{実在する}残余対角のperfect性、
すなわち$\Gamma_R\in\Perf(\Gamma_R^{\mathrm{op}}\otL_R\Gamma_R)$は別の問題である。
所定の量子化のgeneric復元の最優先課題は依然としてこの残余対角のperfect性である。
本節はBFKの既知の公式の具体的な帰結と、その本研究における障害の切分けを与える。
\end{remark}


\section{generic可換復元への正確な還元}\label{sec:reconstruction}
\begin{theorem}[局所saturationからgeneric復元へ]\label{thm:generic}
$S$を滑らかな既約有限型複素スキーム、$s\in S(\CC)$、$R=\OO_{S,s}$とする。
$R$上に仮定\ref{hyp:algebra}の量子化があり、残余対角条件\eqref{eq:SM}が成立するとする。
$K=\CC(S)$と置くと、ある有限分離拡大$L/K$、滑らかな三次超曲面$Y_L$、
rotationを保つenhanced同値
\[
 (\cK_{B_L}^{\mathrm{coh}},\LL_{\OO}(1))
       \simeq (\Ku(Y_L),\rho_{Y_L})
 \tag{CR}\label{eq:reconstruction}
\]
が存在する。ここで左辺は全coherent残余圏である。
\end{theorem}
\begin{proof}
局所proper模型と定理\ref{thm:rotation}から$(\Perf(\Gamma_R),\rho)$を得る。
\eqref{eq:SM}によりsmooth properとなるので、命題\ref{prop:spread}から$s$を含む開集合$U\subset S$へ
pairを降下する。その$s$-ファイバーは元の$(\Ku(X),\rho_X)$である。
定理\ref{er:thm:ER}を適用してpointed étale近傍$V\to U$と$Y/V$を得る。
指定点を通る既約成分を選ぶと、その像は$U$の非空開集合を含む。
étale射は準有限なので、その関数体$L$は$K$の有限分離拡大である。

線束Homの自然な分数体底変更と有限mutationにより、
$\Gamma_R\otL_RL$は実際の$B_L$の線束生成残余圏を表す。
smoothnessは底変更で保存されるので、第\ref{sec:fullness}節を$L$上で適用して
線束生成部分と全coherent残余圏を同定する。
従ってgenericファイバーで\eqref{eq:reconstruction}を得る。
\end{proof}

\begin{corollary}[Poisson成分に対する適用範囲]\label{cor:component}
Poisson既約成分$Z$を支配する滑らかな既約chart $S$と可換点$s$があり、
$\CC(S)/\CC(Z)$が有限分離であるとする。
$s$の局所環上で仮定\ref{hyp:algebra}と\eqref{eq:SM}を満たせば、
$Z$の関数体の有限分離拡大上に\eqref{eq:reconstruction}の可換模型が存在する。
\end{corollary}
\begin{proof}
定理\ref{thm:generic}を適用し、二つの有限分離拡大を合成する。
\end{proof}

\begin{remark}[残る仮定の量化]
各Poisson成分に一つ、この局所入力を実現すればよい。
全成分での量子化自体が本稿から生じるわけではない。
定理\ref{er:thm:ER}は自己同値一つを指定したpairについての定理である。
可換模型から周期性同型を移すことはできるが、特定のmatrix factorization表示の
周期性とのcoherenceまで自動的に同定したとは扱わない。
\end{remark}

\section{非楕円型への適用と未解決の核心}\label{sec:remaining}
親プロジェクトの二つの非楕円型成分では、可換点の局所環は
$R=\CC[r,s]_{(r,s)}$、$N=7$、$n=5$、$a=4$である。
その全次数flatなambientと中心三次商の構成を入力すれば、系\ref{cor:local-rotation}により
\[
 \rho^{-1}=(-1)\RR_{\OO}
\]
が局所残余圏上に存在する。節点型・尖点型の純三次ブロックの滑らかさは使わない。
具体的な係数と中心式は研究ノートに記す。

この二成分に対する残る核心は、$\Gamma_R$の対角が$\Gamma_R^e$上perfectであることにある。
これが得られれば、全coherent比較、pairのspread、定理\ref{er:thm:ER}により
一般点のrotation付き可換模型が有限分離拡大上に存在する。
本稿はこの対角perfect性を証明したものではない。

\subsection{非楕円二成分に対する具体的な可換変形slice}
変数の順を$(x,y,z,a,b,c,u)$とし、
\[
 F_\epsilon=f_\epsilon+ux^2+a^3+b^3+c^3+u(a^2+b^2+c^2)+u^3,
 \quad f_C=xz^2+y^3,\quad f_N=xyz+y^3+z^3
\]
と置く。これらの滑らかさは研究ノートの勾配消去で確認されている。
\begin{proposition}[35係数の横断slice：厳密線形代数の証明書]\label{er:prop:slice}
次の単項式集合はそれぞれ$J(F_C)_3$と$J(F_N)_3$の基底を与える。
\begin{align*}
\mathcal M_C=\{&axy,bxy,cxy,abx,acx,bcx,yz^2,ayz,byz,cyz,uyz,\\
 &aby,acy,auy,bcy,buy,cuy,u^2y,az^2,bz^2,cz^2,uz^2,\\
 &abz,acz,bcz,buz,cuz,u^2z,abc,abu,acu,bcu,bu^2,cu^2,u^3\},\\
\mathcal M_N=\{&abx,acx,bcx,ay^2,by^2,cy^2,ayz,byz,cyz,uyz,\\
 &aby,acy,bcy,buy,cuy,u^2y,az^2,bz^2,cz^2,uz^2,abz,acz,\\
 &auz,bcz,buz,cuz,u^2z,abc,abu,acu,au^2,bcu,bu^2,cu^2,u^3\}.
\end{align*}
従って$V(F_\epsilon+\sum_{m\in\mathcal M_\epsilon}t_m m)$の滑らかな開部分は、
原点で定理\ref{er:thm:criterion}の幾何的接写像が同型となる35次元のsource族である。
\end{proposition}
\begin{proof}
84個の三次単項式を、変数番号の非減少三つ組の辞書式順で並べる。
84行49列の整数行列$M_\epsilon$の列を$z_j\partial_iF_\epsilon$の係数ベクトルとする。
列は$i$を先、$j$を後に増加させて並べる。
上記$\mathcal M_\epsilon$に対応する35行を除いた49行の小行列を$D_\epsilon$とすると、
整数のfraction-free消去により
\[
 \det D_C=2^{15}3^{26},\qquad \det D_N=-2^6 3^{25}
\]
を得る。独立な二つの決定的アルゴリズムによる検算を
\texttt{computations/verify\_etale\_inputs.py}に収める。
これは有限体の標本検査ではなく、上で完全に指定した整数行列の行列式計算である。
残りの行を$U$、採用した行を$P$とすると、
$v\mapsto v_U-(M_\epsilon)_U D_\epsilon^{-1}v_P$は$M_\epsilon$を零へ送り、
$\mathcal M_\epsilon$の単項式上では恒等となる。この二恒等式も有理数上で検算する。
よって84次元の三次形式空間は、49次元の座標変換方向と表示した35次元の空間の直和である。
Euler等式により$F_\epsilon$自身も前者に入る。補題\ref{er:lem:cubic}で接写像の主張を得る。
\end{proof}

\begin{remark}
この基底は局所実現のsourceを具体化するものであり、量子化パラメータ$(r,s)$から
$t_m$への閉公式を計算したものではない。各小拡大の実際の障害類が得られれば、
その消去をこの35係数で行える。smoothness証明書を代替する計算でもない。
\end{remark}

\begin{proposition}[Serre自己同値はsmoothnessを含意しない]\label{prop:warning}
体$k$上の$H=k[\epsilon]/(\epsilon^2)$に対し、$\Perf(H)$はproperでSerre関手が恒等であるがsmoothではない。
\end{proposition}
\begin{proof}
$\epsilon$の係数を取るtraceは非退化な対称Frobenius pairingを与え、$H^\vee\simeq H$は双加群同型である。
従ってperfect加群上のSerre関手は$-\otL_HH^\vee\simeq\mathrm{id}$である。
一方、単純加群$k$は
$\cdots\xrightarrow{\epsilon}H\xrightarrow{\epsilon}H\to k\to0$
という無限最小分解を持ち、$H$-perfectでない。
$H$がsmoothなら補題\ref{lem:pseudoperfect}に反する。
\end{proof}

この例は局所量子化が非smoothであることを示す反例ではない。
今回構成したSerre関手・逆rotationだけでは対角perfectnessの証明にならないことを示している。
対角の持上げ、形式的Artin切断、底上perfectなHom複体という三種類の有限性を区別する必要がある。

\begin{thebibliography}{99}
\bibitem{BS} Jesse Burke and Greg Stevenson,
\emph{The derived category of a graded Gorenstein ring},
arXiv:1507.00830v1 (2015).
\url{https://arxiv.org/abs/1507.00830}.
参照：Definition 5.1, Proposition 4.18, Lemmas 6.1--6.2, Theorem 6.4.
\bibitem{Toen} Bertrand Toën,
\emph{Anneaux de définition des dg-algèbres propres et lisses},
arXiv:math/0611546v3 (2010).
\url{https://arxiv.org/abs/math/0611546}.
参照：Théorème 1, F4, F6.
\bibitem{TV} Bertrand Toën and Michel Vaquié,
\emph{Moduli of objects in dg-categories},
arXiv:math/0503269v5 (2007).
\url{https://arxiv.org/abs/math/0503269}.
参照：Lemmas 2.8, 2.10, Corollary 2.13, Theorem 3.6, Corollaries 3.15, 3.17, 3.19.
\bibitem{Orlov} Dmitri Orlov,
\emph{Smooth and proper noncommutative schemes and gluing of DG categories},
arXiv:1402.7364v5 (2015).
\url{https://arxiv.org/abs/1402.7364}.
参照：Proposition 3.8, Theorem 3.25（体上）.
\bibitem{Repo} nc-cubic project,
\emph{一般次元の局所代数的模型、非楕円型量子化、およびrotation付き可換復元},
commit \texttt{4f2eb328a2249a173f0c658cdb744277f041ff1d}, 2026年11月10日JST.
\url{https://github.com/uedakazushi/nc-cubic}.
一次原稿：\texttt{papers/sections/local\_projective\_model.tex},
\texttt{nonelliptic\_quantization.tex}, \texttt{rotation\_models.tex}.
局所模型の入力、親原稿の形式的復元と、本稿の双加群による局所実現証明を区別する。
\bibitem{ERLowen} Wendy Lowen,
\emph{Hochschild cohomology, the characteristic morphism and derived deformations},
Compositio Mathematica \textbf{144} (2008), 1557--1580.
\url{https://arxiv.org/abs/0707.2602}. 参照：v1, Theorem 4.8, \S4.4.
\bibitem{BFKER} Matthew Ballard, David Favero and Ludmil Katzarkov,
\emph{A category of kernels for equivariant factorizations and its implications for Hodge theory},
Publications Mathématiques de l'IHÉS \textbf{120} (2014), 1--111.
\url{https://doi.org/10.1007/s10240-013-0059-9}.
参照：Theorem 5.39, pp.~82--83.
\bibitem{HRER} Daniel Huybrechts and Jørgen Vold Rennemo,
\emph{Hochschild cohomology versus the Jacobian ring and the Torelli theorem for cubic fourfolds},
Algebraic Geometry \textbf{6} (2019), 76--99.
\url{https://doi.org/10.14231/AG-2019-005}.
参照：Lemma 2.19, \S2.5, \S3.3（cubicでは$n\ge4$）.
\bibitem{JLLZER} Augustinas Jacovskis, Xun Lin, Zhiyu Liu and Shizhuo Zhang,
\emph{Infinitesimal categorical Torelli theorems for Fano threefolds},
arXiv:2203.08187v3 (2023).
\url{https://arxiv.org/abs/2203.08187}. 参照：Theorem 4.4(2).
\bibitem{HuyER} Daniel Huybrechts, \emph{The K3 category of a cubic fourfold},
Compositio Mathematica \textbf{153} (2017), 586--620.
\url{https://doi.org/10.1112/S0010437X16008137}. 参照：Proposition 3.12, \S3.5.
\bibitem{StacksER} The Stacks Project Authors, \emph{The Stacks Project}.
\url{https://stacks.math.columbia.edu/tag/02HX};
\url{https://stacks.math.columbia.edu/tag/055U}.
参照日：2026年11月10日。smoothnessの無限小判定、smooth射のétale局所section。
\bibitem{ERLinZhang} Xun Lin and Shizhuo Zhang,
\emph{Categorical Torelli and Fourier--Mukai partners of hypersurfaces},
arXiv:2611.05915v1 (2026年11月4日).
\url{https://arxiv.org/abs/2611.05915}.
参照：\S2.1, Theorem 2.2, Proposition 2.4.
幾何的な両側projectorを仮定するため、非幾何的targetには直接適用しない。

\end{thebibliography}
\end{document}
